%% LyX 2.4.4 created this file.  For more info, see https://www.lyx.org/.
%% Do not edit unless you really know what you are doing.
\documentclass[english]{article}
\usepackage[T1]{fontenc}
\usepackage[latin9]{inputenc}
\usepackage{enumitem}
\usepackage{amsmath}
\usepackage{amsthm}
\usepackage{amssymb}
\usepackage{geometry}
\geometry{verbose,tmargin=1in,bmargin=1in,lmargin=1in,rmargin=1in}
\usepackage{wasysym}

\makeatletter
%%%%%%%%%%%%%%%%%%%%%%%%%%%%%% Textclass specific LaTeX commands.
\numberwithin{equation}{section}
\numberwithin{figure}{section}
\newlength{\lyxlabelwidth}      % auxiliary length 
\theoremstyle{plain}
\newtheorem*{thm*}{\protect\theoremname}
\theoremstyle{definition}
\newtheorem*{defn*}{\protect\definitionname}

%%%%%%%%%%%%%%%%%%%%%%%%%%%%%% User specified LaTeX commands.
\usepackage{fancyhdr}
\usepackage{lastpage}
\pagestyle{fancy}
\usepackage{enumitem}
\usepackage{ifthen}
\everymath=\expandafter{\the\everymath\displaystyle}
%\DeclareMathSizes{10}{10}{10}{10}
\usepackage{multicol}

\usepackage{tikz}
\usetikzlibrary{patterns}
\usetikzlibrary{plotmarks}
\usepackage{pgfplots}
\pgfplotsset{compat=newest}

\usepackage{calc}
\usepackage[nomessages]{fp}

\usepackage{datetime}
% Added by lyx2lyx
\setlength{\parskip}{\smallskipamount}
\setlength{\parindent}{0pt}

\makeatother

\usepackage{babel}
\providecommand{\definitionname}{Definition}
\providecommand{\theoremname}{Theorem}

\begin{document}
\renewcommand{\author}{Dr.\,James Ashe}

\newcommand{\classNum}{335}

\newcommand{\classSec}{A}

\newcommand{\nameType}{Name: \hspace{-4cm}}

\newcommand{\examType}{Homework 4}

\newcommand{\Date}{Due date: \formatdate{15}{4}{2013}} %%\AdvanceDate[12] \today

%\newdate{my_date}{\day}{\month}{\year}%   
%\AdvanceDate[-5] 
%\displaydate{my_date}

%%First page's header%%%%%%%%%%%%%%%%%%%%%%
\fancypagestyle{firststyle} %%header settings for first pagr
{
	\fancyhead[L]{MTH \classNum{}(\classSec{}) \author{}\\
	\textbf{\examType{}} \Date{}}
	\fancyhead[C]{\textbf{\nameType}}
	\setlength{\headsep}{14pt}
}
%%%%%%%%%%%%%%%%%%%%%%%%%%%%%%%%%%%%%%%%%%

%%Next page's header%%%%%%%%%%%%%%%%%%%%%%%
%\setlist{nolistsep} %eliminates space between list items
\lhead{MTH \classNum{}(\classSec{})} 
\chead{\textbf{\nameType}} 
\cfoot{\thepage{}~of~\pageref{LastPage}} 
\setlength{\headsep}{5pt} 
%%%%%%%%%%%%%%%%%%%%%%%%%%%%%%%%%%%%%%%%%%%

%%Various settings; see the preamble for other settings%%%%%
%%\setenumerate[0]{leftmargin=12pt,itemindent=0pt,labelwidth=0pt}
\renewcommand{\labelenumi}{\textbf{\arabic{enumi}.}}
\renewcommand{\labelenumii}{\textbf{\alph{enumii}.}}
\thispagestyle{firststyle}
%%%%%%%%%%%%%%%%%%%%%%%%%%%%%%%%%%%%%%%%%%

\textbf{Homework Problems. }\emph{Solutions should be written/typed
in numerical order and clearly labeled. Unsolved problems should be
included but left blank.}
\begin{enumerate}[resume]
\item Let $a,b,c\in\mathbb{Z}$. Prove the following:

\begin{enumerate}[resume]
\item If $a|b$ and $b|c$ then $a|c$.
\item If $a|b$ and $a|c$ then $a|b\pm c$.\vfill
\end{enumerate}
\item Show that the set of integers, $\mathbb{Z}$, is a group under addition.
\vfill
\item Prove that the set of even integers is a group under addition. \vfill
\item Show that the following are \emph{not} groups under the given operation.

\begin{enumerate}
\item $\mathbb{Z}$ under multiplication.
\item $S=\left\{ z|z\in\mathbb{Z}\mbox{ and }z=\sqrt{5}\right\} $under
addition.
\item $\mathbb{Q}$, the set of rationals under multiplication.
\item The set of odd integers under addition.\vfill
\end{enumerate}
\item Show that $\left\{ 1,2,3\right\} $ under multiplication modulo $4$
is \emph{not }a group but that $\left\{ 1,2,3,4\right\} $ under multiplication
modulo $5$ is a group. (Example: $2\cdot3=6$ and $6=2\bmod4$ so
$[6]=[2]$ and thus $2\cdot3$ \emph{is} in $\left\{ 1,2,3\right\} $).
\vfill
\item Prove that the set of even integers is a group under addition. (I
mistakenly repeated problem 3. For a more challenging problem you
might try showing that integers of $n$ multiples for $n\in\mathbb{Z}^{+}$
are a group under addition)\vfill
\item Let $S=\left\{ a+b\sqrt{5}|\,a,b\mbox{ are not both zero and }a,b\in\mathbb{Q}\right\} $.
Prove that $S$ is a group. (Hint: rationalize the denominator to
show that the inverse of $a+b\sqrt{5}$ is in $S$.)\vfill
\item Prove that the set of all $2\times2$ matrices with determinant $1$
and entries from $\mathbb{R}$ is a group under matrix multiplication.
\vfill
\item Suppose that if $G$ is a group such that if $a,b,c\in G$ and $ab=ca$
then $b=c$. Prove that $G$ is an Abelian group, i.e., that $ab=ba$
for all $a,b\in G$. Hint: consider the element $aba=aba$.\vfill
\item Prove that $G$ is an Abelian group if and only is $\left(ab\right)^{-1}=a^{-1}b^{-1}$
for all $a,b\in G$. \vfill\newpage
\end{enumerate}
\textbf{Presentation problems.}
\begin{enumerate}
\item Suppose that $a\neq0$, $b\neq0$, and $\mbox{lcm}(a,b)=m$. Prove
that if $n$ is a common multiple of $a$ and $b$ then $m|n$. 
\end{enumerate}
\begin{enumerate}[resume]
\item Let $a\neq0$, $b\neq0$, and $d=\gcd\left(a,b\right)$. Prove that
if $d'$ is any common divisor of $a$ and $b$ then $d'|d$.
\end{enumerate}
\begin{thm*}
\textbf{\emph{The Fundamental Theorem of Arithmetic. }}Every integer
greater than 1 is either a prime or the unique product of primes. 
\end{thm*}
\begin{enumerate}[resume]
\item If $c^{2}=ab$ and $\gcd\left(a,b\right)=1$, prove that $a$ and
$b$ are perfect squares (Hint: By the theorem above, $a$ and $b$
can be written as the product of primes. They are perfect squares
iff. each prime in this prime factorization has an even power). 
\item Prove that if in a group $G$ that $\left(ab\right)^{2}=a^{2}b^{2}$
then $G$ is Abelian. 
\item Consider a square placed on Cartesian plane such that the square's
center is at the origin and its corners are on the axes (like a diamond;
see the figure below). Any motion of the square which keeps the center
at the origin and corners on the axes has the same result as one of
eight motions: $r_{0}\equiv\mbox{ rotation of }0^{\circ}$, $r_{1}\equiv\mbox{ rotation of }90^{\circ}$,
$r_{3}\equiv\mbox{ rotation of }180^{\circ}$, $d_{x}\equiv\mbox{ }x-\mbox{axis reflection}$,
$h_{+}\equiv\,y=x\mbox{ reflection}$, or $h_{-}\equiv y=-x\mbox{ reflection}$.
\\
\hspace*{.5cm}\newcommand{\xmax}{1.5}
\newcommand{\xmin}{-1.5}
\newcommand{\xstep}{0} %must be xmin+n
\newcommand{\ymax}{1.5}
\newcommand{\ymin}{-1.5}
\newcommand{\ystep}{0} %must be ymin+n
%%%%%%%
\begin{tikzpicture} 
\begin{axis}[height=3.5cm, 
	width=3.5cm, 
	axis lines=middle,
	x axis line style={-},y axis line style={-}, 
	grid=none,
	%axis line style={-latex}, 
	xmin=\xmin,xmax=\xmax, 
	ymin=\ymin,ymax=\ymax,
	%restrict y to domain=-1:10,
	no markers, 
	xlabel={$$},ylabel={$$},
	ticks=none,
	]
	
	\addplot[very thick,blue] coordinates{(1,0) (0,1)};
	\addplot[very thick,blue] coordinates{(0,1) (-1,0)};
	\addplot[very thick,blue] coordinates{(-1,0) (0,-1)};
	\addplot[very thick,blue] coordinates{(0,-1) (1,0)};

	\node at (axis cs:1.3,.3) {$1$};
	\node at (axis cs:-.3,1.3) {$2$};
	\node at (axis cs:-1.3,-.3) {$3$};
	\node at (axis cs:.3,-1.3) {$4$};
	
\end{axis}
\end{tikzpicture}\\
Show that $D_{4}=\left\{ r_{0},r_{1},r_{3},d_{x},d_{y},h_{+},h_{-}\right\} $
is a group under composition of these motions. The set $D_{4}$ is
called the \emph{dihedral group of degree 4}. 
\end{enumerate}
\begin{defn*}
Let $a$ and $n$ be integers with $n>0$. The \emph{congruence class
of a modulo n }(denoted $[a]$) is the set of integers that are congruent
to $a$ modulo $n$, i.e., 
\[
[a]=\left\{ b|b\in\mathbb{Z}\mbox{ and }b=a\bmod n\right\} 
\]
\end{defn*}
For example, the congruence class of $3$ modulo $5$ is $[3]=\left\{ \dots,-2,3,8,13,\dots\right\} $.

\vspace{0cm}
\begin{defn*}
The set of all congruence classes modulo $n$ is denoted $\mathbb{Z}_{n}$
(read ``$Z$ mod $n$''). 

\begin{itemize}
\item $\left|\mathbb{Z}_{n}\right|=n$. 
\item Addition in $\mathbb{Z}_{n}$ is defined as $\left[a\right]+\left[b\right]=\left[a+b\right]$.
\item Multiplication in $\mathbb{Z}_{n}$ is defined as $\left[a\right]\cdot\left[b\right]=\left[a\cdot b\right]$.
\end{itemize}
\end{defn*}
For example, $\mathbb{Z}_{4}=\left\{ \left[0\right],\left[1\right],\left[2\right],\left[3\right]\right\} $
for convenience we write $\mathbb{Z}_{4}=\left\{ 0,1,2,3\right\} $.
$2+3=5$ and $5=1\bmod4$ so $\left[5\right]=\left[1\right]$.
\begin{enumerate}[resume]
\item Show that $\mathbb{Z}_{2}\times\mathbb{Z}_{3}$ is a group under modular
addition. 
\end{enumerate}
\begin{defn*}
A subset $H$ of a group $G$ is a \emph{subgroup }of $G$ is $H$
is also a group under the operation in $G$.
\end{defn*}
\begin{enumerate}[resume]
\item Let $\mathbb{Z}_{6}=\left\{ 0,1,2,3,4,5\right\} $ and $\mathbb{Z}_{4}=\left\{ 0,1,2,3\right\} $.
$\mathbb{Z}_{6}\times\mathbb{Z}_{4}$ is a group under modular addition.
Show that $H=\left\{ \left(0,0\right),\left(3,0\right),\left(0,2\right),\left(3,2\right)\right\} $
is a subgroup. 
\item [$\smiley$\textbf{.}]Two sets are said to be \emph{equal }if there
exists a bijective function between them; we write $\left|A\right|=\left|B\right|$.
$B$ is said to have \emph{strictly greater cardinality} than $B$
if there is a one-to-one function from $A$ into $B$, but there does
not exist a bijective function from $A$ into $B$; we write$\left|A\right|<\left|B\right|$.
Prove that $\left|\mathbb{N}\right|<\left|\mathbb{R}\right|$.
\end{enumerate}

\end{document}
